\documentclass[11pt]{article}
\usepackage{fontspec}

    \usepackage[breakable]{tcolorbox}
    \usepackage{parskip} % Stop auto-indenting (to mimic markdown behaviour)
    

    % Basic figure setup, for now with no caption control since it's done
    % automatically by Pandoc (which extracts ![](path) syntax from Markdown).
    \usepackage{graphicx}
    % Keep aspect ratio if custom image width or height is specified
    \setkeys{Gin}{keepaspectratio}
    % Maintain compatibility with old templates. Remove in nbconvert 6.0
    \let\Oldincludegraphics\includegraphics
    % Ensure that by default, figures have no caption (until we provide a
    % proper Figure object with a Caption API and a way to capture that
    % in the conversion process - todo).
    \usepackage{caption}

    \usepackage{float}
    \floatplacement{figure}{H} % forces figures to be placed at the correct location
    \usepackage{xcolor} % Allow colors to be defined
    \usepackage{enumerate} % Needed for markdown enumerations to work
    \usepackage{geometry} % Used to adjust the document margins
    \usepackage{amsmath} % Equations
    \usepackage{amssymb} % Equations
    \usepackage{textcomp} % defines textquotesingle
    % Hack from http://tex.stackexchange.com/a/47451/13684:
    \AtBeginDocument{%
        \def\PYZsq{\textquotesingle}% Upright quotes in Pygmentized code
    }
    \usepackage{upquote} % Upright quotes for verbatim code
    \usepackage{eurosym} % defines \euro

    \usepackage{iftex}
    \ifPDFTeX
        
        \IfFileExists{alphabeta.sty}{
              \usepackage{alphabeta}
          }{
              \usepackage[mathletters]{ucs}
              
          }
    \else
        \usepackage{fontspec}
        \usepackage{unicode-math}
    \fi

    \usepackage{fancyvrb} % verbatim replacement that allows latex
    \usepackage{grffile} % extends the file name processing of package graphics
                         % to support a larger range
    \makeatletter % fix for old versions of grffile with XeLaTeX
    \@ifpackagelater{grffile}{2019/11/01}
    {
      % Do nothing on new versions
    }
    {
      \def\Gread@@xetex#1{%
        \IfFileExists{"\Gin@base".bb}%
        {\Gread@eps{\Gin@base.bb}}%
        {\Gread@@xetex@aux#1}%
      }
    }
    \makeatother
    \usepackage[Export]{adjustbox} % Used to constrain images to a maximum size
    \adjustboxset{max size={0.9\linewidth}{0.9\paperheight}}

    % The hyperref package gives us a pdf with properly built
    % internal navigation ('pdf bookmarks' for the table of contents,
    % internal cross-reference links, web links for URLs, etc.)
    \usepackage{hyperref}
    % The default LaTeX title has an obnoxious amount of whitespace. By default,
    % titling removes some of it. It also provides customization options.
    \usepackage{titling}
    \usepackage{longtable} % longtable support required by pandoc >1.10
    \usepackage{booktabs}  % table support for pandoc > 1.12.2
    \usepackage{array}     % table support for pandoc >= 2.11.3
    \usepackage{calc}      % table minipage width calculation for pandoc >= 2.11.1
    \usepackage[inline]{enumitem} % IRkernel/repr support (it uses the enumerate* environment)
    \usepackage[normalem]{ulem} % ulem is needed to support strikethroughs (\sout)
                                % normalem makes italics be italics, not underlines
    \usepackage{soul}      % strikethrough (\st) support for pandoc >= 3.0.0
    \usepackage{mathrsfs}
    

    
    % Colors for the hyperref package
    \definecolor{urlcolor}{rgb}{0,.145,.698}
    \definecolor{linkcolor}{rgb}{.71,0.21,0.01}
    \definecolor{citecolor}{rgb}{.12,.54,.11}

    % ANSI colors
    \definecolor{ansi-black}{HTML}{3E424D}
    \definecolor{ansi-black-intense}{HTML}{282C36}
    \definecolor{ansi-red}{HTML}{E75C58}
    \definecolor{ansi-red-intense}{HTML}{B22B31}
    \definecolor{ansi-green}{HTML}{00A250}
    \definecolor{ansi-green-intense}{HTML}{007427}
    \definecolor{ansi-yellow}{HTML}{DDB62B}
    \definecolor{ansi-yellow-intense}{HTML}{B27D12}
    \definecolor{ansi-blue}{HTML}{208FFB}
    \definecolor{ansi-blue-intense}{HTML}{0065CA}
    \definecolor{ansi-magenta}{HTML}{D160C4}
    \definecolor{ansi-magenta-intense}{HTML}{A03196}
    \definecolor{ansi-cyan}{HTML}{60C6C8}
    \definecolor{ansi-cyan-intense}{HTML}{258F8F}
    \definecolor{ansi-white}{HTML}{C5C1B4}
    \definecolor{ansi-white-intense}{HTML}{A1A6B2}
    \definecolor{ansi-default-inverse-fg}{HTML}{FFFFFF}
    \definecolor{ansi-default-inverse-bg}{HTML}{000000}

    % common color for the border for error outputs.
    \definecolor{outerrorbackground}{HTML}{FFDFDF}

    % commands and environments needed by pandoc snippets
    % extracted from the output of `pandoc -s`
    \providecommand{\tightlist}{%
      \setlength{\itemsep}{0pt}\setlength{\parskip}{0pt}}
    \DefineVerbatimEnvironment{Highlighting}{Verbatim}{commandchars=\\\{\}}
    % Add ',fontsize=\small' for more characters per line
    \newenvironment{Shaded}{}{}
    \newcommand{\KeywordTok}[1]{\textcolor[rgb]{0.00,0.44,0.13}{\textbf{{#1}}}}
    \newcommand{\DataTypeTok}[1]{\textcolor[rgb]{0.56,0.13,0.00}{{#1}}}
    \newcommand{\DecValTok}[1]{\textcolor[rgb]{0.25,0.63,0.44}{{#1}}}
    \newcommand{\BaseNTok}[1]{\textcolor[rgb]{0.25,0.63,0.44}{{#1}}}
    \newcommand{\FloatTok}[1]{\textcolor[rgb]{0.25,0.63,0.44}{{#1}}}
    \newcommand{\CharTok}[1]{\textcolor[rgb]{0.25,0.44,0.63}{{#1}}}
    \newcommand{\StringTok}[1]{\textcolor[rgb]{0.25,0.44,0.63}{{#1}}}
    \newcommand{\CommentTok}[1]{\textcolor[rgb]{0.38,0.63,0.69}{\textit{{#1}}}}
    \newcommand{\OtherTok}[1]{\textcolor[rgb]{0.00,0.44,0.13}{{#1}}}
    \newcommand{\AlertTok}[1]{\textcolor[rgb]{1.00,0.00,0.00}{\textbf{{#1}}}}
    \newcommand{\FunctionTok}[1]{\textcolor[rgb]{0.02,0.16,0.49}{{#1}}}
    \newcommand{\RegionMarkerTok}[1]{{#1}}
    \newcommand{\ErrorTok}[1]{\textcolor[rgb]{1.00,0.00,0.00}{\textbf{{#1}}}}
    \newcommand{\NormalTok}[1]{{#1}}

    % Additional commands for more recent versions of Pandoc
    \newcommand{\ConstantTok}[1]{\textcolor[rgb]{0.53,0.00,0.00}{{#1}}}
    \newcommand{\SpecialCharTok}[1]{\textcolor[rgb]{0.25,0.44,0.63}{{#1}}}
    \newcommand{\VerbatimStringTok}[1]{\textcolor[rgb]{0.25,0.44,0.63}{{#1}}}
    \newcommand{\SpecialStringTok}[1]{\textcolor[rgb]{0.73,0.40,0.53}{{#1}}}
    \newcommand{\ImportTok}[1]{{#1}}
    \newcommand{\DocumentationTok}[1]{\textcolor[rgb]{0.73,0.13,0.13}{\textit{{#1}}}}
    \newcommand{\AnnotationTok}[1]{\textcolor[rgb]{0.38,0.63,0.69}{\textbf{\textit{{#1}}}}}
    \newcommand{\CommentVarTok}[1]{\textcolor[rgb]{0.38,0.63,0.69}{\textbf{\textit{{#1}}}}}
    \newcommand{\VariableTok}[1]{\textcolor[rgb]{0.10,0.09,0.49}{{#1}}}
    \newcommand{\ControlFlowTok}[1]{\textcolor[rgb]{0.00,0.44,0.13}{\textbf{{#1}}}}
    \newcommand{\OperatorTok}[1]{\textcolor[rgb]{0.40,0.40,0.40}{{#1}}}
    \newcommand{\BuiltInTok}[1]{{#1}}
    \newcommand{\ExtensionTok}[1]{{#1}}
    \newcommand{\PreprocessorTok}[1]{\textcolor[rgb]{0.74,0.48,0.00}{{#1}}}
    \newcommand{\AttributeTok}[1]{\textcolor[rgb]{0.49,0.56,0.16}{{#1}}}
    \newcommand{\InformationTok}[1]{\textcolor[rgb]{0.38,0.63,0.69}{\textbf{\textit{{#1}}}}}
    \newcommand{\WarningTok}[1]{\textcolor[rgb]{0.38,0.63,0.69}{\textbf{\textit{{#1}}}}}
    \makeatletter
    \newsavebox\pandoc@box
    \newcommand*\pandocbounded[1]{%
      \sbox\pandoc@box{#1}%
      % scaling factors for width and height
      \Gscale@div\@tempa\textheight{\dimexpr\ht\pandoc@box+\dp\pandoc@box\relax}%
      \Gscale@div\@tempb\linewidth{\wd\pandoc@box}%
      % select the smaller of both
      \ifdim\@tempb\p@<\@tempa\p@
        \let\@tempa\@tempb
      \fi
      % scaling accordingly (\@tempa < 1)
      \ifdim\@tempa\p@<\p@
        \scalebox{\@tempa}{\usebox\pandoc@box}%
      % scaling not needed, use as it is
      \else
        \usebox{\pandoc@box}%
      \fi
    }
    \makeatother

    % Define a nice break command that doesn't care if a line doesn't already
    % exist.
    \def\br{\hspace*{\fill} \\* }
    % Math Jax compatibility definitions
    \def\gt{>}
    \def\lt{<}
    \let\Oldtex\TeX
    \let\Oldlatex\LaTeX
    \renewcommand{\TeX}{\textrm{\Oldtex}}
    \renewcommand{\LaTeX}{\textrm{\Oldlatex}}
    % Document parameters
    % Document title
    
    
    
    
    
    
    
    
% Pygments definitions
\makeatletter
\def\PY@reset{\let\PY@it=\relax \let\PY@bf=\relax%
    \let\PY@ul=\relax \let\PY@tc=\relax%
    \let\PY@bc=\relax \let\PY@ff=\relax}
\def\PY@tok#1{\csname PY@tok@#1\endcsname}
\def\PY@toks#1+{\ifx\relax#1\empty\else%
    \PY@tok{#1}\expandafter\PY@toks\fi}
\def\PY@do#1{\PY@bc{\PY@tc{\PY@ul{%
    \PY@it{\PY@bf{\PY@ff{#1}}}}}}}
\def\PY#1#2{\PY@reset\PY@toks#1+\relax+\PY@do{#2}}

\@namedef{PY@tok@w}{\def\PY@tc##1{\textcolor[rgb]{0.73,0.73,0.73}{##1}}}
\@namedef{PY@tok@c}{\let\PY@it=\textit\def\PY@tc##1{\textcolor[rgb]{0.24,0.48,0.48}{##1}}}
\@namedef{PY@tok@cp}{\def\PY@tc##1{\textcolor[rgb]{0.61,0.40,0.00}{##1}}}
\@namedef{PY@tok@k}{\let\PY@bf=\textbf\def\PY@tc##1{\textcolor[rgb]{0.00,0.50,0.00}{##1}}}
\@namedef{PY@tok@kp}{\def\PY@tc##1{\textcolor[rgb]{0.00,0.50,0.00}{##1}}}
\@namedef{PY@tok@kt}{\def\PY@tc##1{\textcolor[rgb]{0.69,0.00,0.25}{##1}}}
\@namedef{PY@tok@o}{\def\PY@tc##1{\textcolor[rgb]{0.40,0.40,0.40}{##1}}}
\@namedef{PY@tok@ow}{\let\PY@bf=\textbf\def\PY@tc##1{\textcolor[rgb]{0.67,0.13,1.00}{##1}}}
\@namedef{PY@tok@nb}{\def\PY@tc##1{\textcolor[rgb]{0.00,0.50,0.00}{##1}}}
\@namedef{PY@tok@nf}{\def\PY@tc##1{\textcolor[rgb]{0.00,0.00,1.00}{##1}}}
\@namedef{PY@tok@nc}{\let\PY@bf=\textbf\def\PY@tc##1{\textcolor[rgb]{0.00,0.00,1.00}{##1}}}
\@namedef{PY@tok@nn}{\let\PY@bf=\textbf\def\PY@tc##1{\textcolor[rgb]{0.00,0.00,1.00}{##1}}}
\@namedef{PY@tok@ne}{\let\PY@bf=\textbf\def\PY@tc##1{\textcolor[rgb]{0.80,0.25,0.22}{##1}}}
\@namedef{PY@tok@nv}{\def\PY@tc##1{\textcolor[rgb]{0.10,0.09,0.49}{##1}}}
\@namedef{PY@tok@no}{\def\PY@tc##1{\textcolor[rgb]{0.53,0.00,0.00}{##1}}}
\@namedef{PY@tok@nl}{\def\PY@tc##1{\textcolor[rgb]{0.46,0.46,0.00}{##1}}}
\@namedef{PY@tok@ni}{\let\PY@bf=\textbf\def\PY@tc##1{\textcolor[rgb]{0.44,0.44,0.44}{##1}}}
\@namedef{PY@tok@na}{\def\PY@tc##1{\textcolor[rgb]{0.41,0.47,0.13}{##1}}}
\@namedef{PY@tok@nt}{\let\PY@bf=\textbf\def\PY@tc##1{\textcolor[rgb]{0.00,0.50,0.00}{##1}}}
\@namedef{PY@tok@nd}{\def\PY@tc##1{\textcolor[rgb]{0.67,0.13,1.00}{##1}}}
\@namedef{PY@tok@s}{\def\PY@tc##1{\textcolor[rgb]{0.73,0.13,0.13}{##1}}}
\@namedef{PY@tok@sd}{\let\PY@it=\textit\def\PY@tc##1{\textcolor[rgb]{0.73,0.13,0.13}{##1}}}
\@namedef{PY@tok@si}{\let\PY@bf=\textbf\def\PY@tc##1{\textcolor[rgb]{0.64,0.35,0.47}{##1}}}
\@namedef{PY@tok@se}{\let\PY@bf=\textbf\def\PY@tc##1{\textcolor[rgb]{0.67,0.36,0.12}{##1}}}
\@namedef{PY@tok@sr}{\def\PY@tc##1{\textcolor[rgb]{0.64,0.35,0.47}{##1}}}
\@namedef{PY@tok@ss}{\def\PY@tc##1{\textcolor[rgb]{0.10,0.09,0.49}{##1}}}
\@namedef{PY@tok@sx}{\def\PY@tc##1{\textcolor[rgb]{0.00,0.50,0.00}{##1}}}
\@namedef{PY@tok@m}{\def\PY@tc##1{\textcolor[rgb]{0.40,0.40,0.40}{##1}}}
\@namedef{PY@tok@gh}{\let\PY@bf=\textbf\def\PY@tc##1{\textcolor[rgb]{0.00,0.00,0.50}{##1}}}
\@namedef{PY@tok@gu}{\let\PY@bf=\textbf\def\PY@tc##1{\textcolor[rgb]{0.50,0.00,0.50}{##1}}}
\@namedef{PY@tok@gd}{\def\PY@tc##1{\textcolor[rgb]{0.63,0.00,0.00}{##1}}}
\@namedef{PY@tok@gi}{\def\PY@tc##1{\textcolor[rgb]{0.00,0.52,0.00}{##1}}}
\@namedef{PY@tok@gr}{\def\PY@tc##1{\textcolor[rgb]{0.89,0.00,0.00}{##1}}}
\@namedef{PY@tok@ge}{\let\PY@it=\textit}
\@namedef{PY@tok@gs}{\let\PY@bf=\textbf}
\@namedef{PY@tok@ges}{\let\PY@bf=\textbf\let\PY@it=\textit}
\@namedef{PY@tok@gp}{\let\PY@bf=\textbf\def\PY@tc##1{\textcolor[rgb]{0.00,0.00,0.50}{##1}}}
\@namedef{PY@tok@go}{\def\PY@tc##1{\textcolor[rgb]{0.44,0.44,0.44}{##1}}}
\@namedef{PY@tok@gt}{\def\PY@tc##1{\textcolor[rgb]{0.00,0.27,0.87}{##1}}}
\@namedef{PY@tok@err}{\def\PY@bc##1{{\setlength{\fboxsep}{\string -\fboxrule}\fcolorbox[rgb]{1.00,0.00,0.00}{1,1,1}{\strut ##1}}}}
\@namedef{PY@tok@kc}{\let\PY@bf=\textbf\def\PY@tc##1{\textcolor[rgb]{0.00,0.50,0.00}{##1}}}
\@namedef{PY@tok@kd}{\let\PY@bf=\textbf\def\PY@tc##1{\textcolor[rgb]{0.00,0.50,0.00}{##1}}}
\@namedef{PY@tok@kn}{\let\PY@bf=\textbf\def\PY@tc##1{\textcolor[rgb]{0.00,0.50,0.00}{##1}}}
\@namedef{PY@tok@kr}{\let\PY@bf=\textbf\def\PY@tc##1{\textcolor[rgb]{0.00,0.50,0.00}{##1}}}
\@namedef{PY@tok@bp}{\def\PY@tc##1{\textcolor[rgb]{0.00,0.50,0.00}{##1}}}
\@namedef{PY@tok@fm}{\def\PY@tc##1{\textcolor[rgb]{0.00,0.00,1.00}{##1}}}
\@namedef{PY@tok@vc}{\def\PY@tc##1{\textcolor[rgb]{0.10,0.09,0.49}{##1}}}
\@namedef{PY@tok@vg}{\def\PY@tc##1{\textcolor[rgb]{0.10,0.09,0.49}{##1}}}
\@namedef{PY@tok@vi}{\def\PY@tc##1{\textcolor[rgb]{0.10,0.09,0.49}{##1}}}
\@namedef{PY@tok@vm}{\def\PY@tc##1{\textcolor[rgb]{0.10,0.09,0.49}{##1}}}
\@namedef{PY@tok@sa}{\def\PY@tc##1{\textcolor[rgb]{0.73,0.13,0.13}{##1}}}
\@namedef{PY@tok@sb}{\def\PY@tc##1{\textcolor[rgb]{0.73,0.13,0.13}{##1}}}
\@namedef{PY@tok@sc}{\def\PY@tc##1{\textcolor[rgb]{0.73,0.13,0.13}{##1}}}
\@namedef{PY@tok@dl}{\def\PY@tc##1{\textcolor[rgb]{0.73,0.13,0.13}{##1}}}
\@namedef{PY@tok@s2}{\def\PY@tc##1{\textcolor[rgb]{0.73,0.13,0.13}{##1}}}
\@namedef{PY@tok@sh}{\def\PY@tc##1{\textcolor[rgb]{0.73,0.13,0.13}{##1}}}
\@namedef{PY@tok@s1}{\def\PY@tc##1{\textcolor[rgb]{0.73,0.13,0.13}{##1}}}
\@namedef{PY@tok@mb}{\def\PY@tc##1{\textcolor[rgb]{0.40,0.40,0.40}{##1}}}
\@namedef{PY@tok@mf}{\def\PY@tc##1{\textcolor[rgb]{0.40,0.40,0.40}{##1}}}
\@namedef{PY@tok@mh}{\def\PY@tc##1{\textcolor[rgb]{0.40,0.40,0.40}{##1}}}
\@namedef{PY@tok@mi}{\def\PY@tc##1{\textcolor[rgb]{0.40,0.40,0.40}{##1}}}
\@namedef{PY@tok@il}{\def\PY@tc##1{\textcolor[rgb]{0.40,0.40,0.40}{##1}}}
\@namedef{PY@tok@mo}{\def\PY@tc##1{\textcolor[rgb]{0.40,0.40,0.40}{##1}}}
\@namedef{PY@tok@ch}{\let\PY@it=\textit\def\PY@tc##1{\textcolor[rgb]{0.24,0.48,0.48}{##1}}}
\@namedef{PY@tok@cm}{\let\PY@it=\textit\def\PY@tc##1{\textcolor[rgb]{0.24,0.48,0.48}{##1}}}
\@namedef{PY@tok@cpf}{\let\PY@it=\textit\def\PY@tc##1{\textcolor[rgb]{0.24,0.48,0.48}{##1}}}
\@namedef{PY@tok@c1}{\let\PY@it=\textit\def\PY@tc##1{\textcolor[rgb]{0.24,0.48,0.48}{##1}}}
\@namedef{PY@tok@cs}{\let\PY@it=\textit\def\PY@tc##1{\textcolor[rgb]{0.24,0.48,0.48}{##1}}}

\def\PYZbs{\char`\\}
\def\PYZus{\char`\_}
\def\PYZob{\char`\{}
\def\PYZcb{\char`\}}
\def\PYZca{\char`\^}
\def\PYZam{\char`\&}
\def\PYZlt{\char`\<}
\def\PYZgt{\char`\>}
\def\PYZsh{\char`\#}
\def\PYZpc{\char`\%}
\def\PYZdl{\char`\$}
\def\PYZhy{\char`\-}
\def\PYZsq{\char`\'}
\def\PYZdq{\char`\"}
\def\PYZti{\char`\~}
% for compatibility with earlier versions
\def\PYZat{@}
\def\PYZlb{[}
\def\PYZrb{]}
\makeatother


    % For linebreaks inside Verbatim environment from package fancyvrb.
    \makeatletter
        \newbox\Wrappedcontinuationbox
        \newbox\Wrappedvisiblespacebox
        \newcommand*\Wrappedvisiblespace {\textcolor{red}{\textvisiblespace}}
        \newcommand*\Wrappedcontinuationsymbol {\textcolor{red}{\llap{\tiny$\m@th\hookrightarrow$}}}
        \newcommand*\Wrappedcontinuationindent {3ex }
        \newcommand*\Wrappedafterbreak {\kern\Wrappedcontinuationindent\copy\Wrappedcontinuationbox}
        % Take advantage of the already applied Pygments mark-up to insert
        % potential linebreaks for TeX processing.
        %        {, <, #, %, $, ' and ": go to next line.
        %        _, }, ^, &, >, - and ~: stay at end of broken line.
        % Use of \textquotesingle for straight quote.
        \newcommand*\Wrappedbreaksatspecials {%
            \def\PYGZus{\discretionary{\char`\_}{\Wrappedafterbreak}{\char`\_}}%
            \def\PYGZob{\discretionary{}{\Wrappedafterbreak\char`\{}{\char`\{}}%
            \def\PYGZcb{\discretionary{\char`\}}{\Wrappedafterbreak}{\char`\}}}%
            \def\PYGZca{\discretionary{\char`\^}{\Wrappedafterbreak}{\char`\^}}%
            \def\PYGZam{\discretionary{\char`\&}{\Wrappedafterbreak}{\char`\&}}%
            \def\PYGZlt{\discretionary{}{\Wrappedafterbreak\char`\<}{\char`\<}}%
            \def\PYGZgt{\discretionary{\char`\>}{\Wrappedafterbreak}{\char`\>}}%
            \def\PYGZsh{\discretionary{}{\Wrappedafterbreak\char`\#}{\char`\#}}%
            \def\PYGZpc{\discretionary{}{\Wrappedafterbreak\char`\%}{\char`\%}}%
            \def\PYGZdl{\discretionary{}{\Wrappedafterbreak\char`\$}{\char`\$}}%
            \def\PYGZhy{\discretionary{\char`\-}{\Wrappedafterbreak}{\char`\-}}%
            \def\PYGZsq{\discretionary{}{\Wrappedafterbreak\textquotesingle}{\textquotesingle}}%
            \def\PYGZdq{\discretionary{}{\Wrappedafterbreak\char`\"}{\char`\"}}%
            \def\PYGZti{\discretionary{\char`\~}{\Wrappedafterbreak}{\char`\~}}%
        }
        % Some characters . , ; ? ! / are not pygmentized.
        % This macro makes them "active" and they will insert potential linebreaks
        \newcommand*\Wrappedbreaksatpunct {%
            \lccode`\~`\.\lowercase{\def~}{\discretionary{\hbox{\char`\.}}{\Wrappedafterbreak}{\hbox{\char`\.}}}%
            \lccode`\~`\,\lowercase{\def~}{\discretionary{\hbox{\char`\,}}{\Wrappedafterbreak}{\hbox{\char`\,}}}%
            \lccode`\~`\;\lowercase{\def~}{\discretionary{\hbox{\char`\;}}{\Wrappedafterbreak}{\hbox{\char`\;}}}%
            \lccode`\~`\:\lowercase{\def~}{\discretionary{\hbox{\char`\:}}{\Wrappedafterbreak}{\hbox{\char`\:}}}%
            \lccode`\~`\?\lowercase{\def~}{\discretionary{\hbox{\char`\?}}{\Wrappedafterbreak}{\hbox{\char`\?}}}%
            \lccode`\~`\!\lowercase{\def~}{\discretionary{\hbox{\char`\!}}{\Wrappedafterbreak}{\hbox{\char`\!}}}%
            \lccode`\~`\/\lowercase{\def~}{\discretionary{\hbox{\char`\/}}{\Wrappedafterbreak}{\hbox{\char`\/}}}%
            \catcode`\.\active
            \catcode`\,\active
            \catcode`\;\active
            \catcode`\:\active
            \catcode`\?\active
            \catcode`\!\active
            \catcode`\/\active
            \lccode`\~`\~
        }
    \makeatother

    \let\OriginalVerbatim=\Verbatim
    \makeatletter
    \renewcommand{\Verbatim}[1][1]{%
        %\parskip\z@skip
        \sbox\Wrappedcontinuationbox {\Wrappedcontinuationsymbol}%
        \sbox\Wrappedvisiblespacebox {\FV@SetupFont\Wrappedvisiblespace}%
        \def\FancyVerbFormatLine ##1{\hsize\linewidth
            \vtop{\raggedright\hyphenpenalty\z@\exhyphenpenalty\z@
                \doublehyphendemerits\z@\finalhyphendemerits\z@
                \strut ##1\strut}%
        }%
        % If the linebreak is at a space, the latter will be displayed as visible
        % space at end of first line, and a continuation symbol starts next line.
        % Stretch/shrink are however usually zero for typewriter font.
        \def\FV@Space {%
            \nobreak\hskip\z@ plus\fontdimen3\font minus\fontdimen4\font
            \discretionary{\copy\Wrappedvisiblespacebox}{\Wrappedafterbreak}
            {\kern\fontdimen2\font}%
        }%

        % Allow breaks at special characters using \PYG... macros.
        \Wrappedbreaksatspecials
        % Breaks at punctuation characters . , ; ? ! and / need catcode=\active
        \OriginalVerbatim[#1,codes*=\Wrappedbreaksatpunct]%
    }
    \makeatother

    % Exact colors from NB
    \definecolor{incolor}{HTML}{303F9F}
    \definecolor{outcolor}{HTML}{D84315}
    \definecolor{cellborder}{HTML}{CFCFCF}
    \definecolor{cellbackground}{HTML}{F7F7F7}

    % prompt
    \makeatletter
    \newcommand{\boxspacing}{\kern\kvtcb@left@rule\kern\kvtcb@boxsep}
    \makeatother
    \newcommand{\prompt}[4]{
        {\ttfamily\llap{{\color{#2}[#3]:\hspace{3pt}#4}}\vspace{-\baselineskip}}
    }
    

    
    % Prevent overflowing lines due to hard-to-break entities
    \sloppy
    % Setup hyperref package
    \hypersetup{
      breaklinks=true,  % so long urls are correctly broken across lines
      colorlinks=true,
      urlcolor=urlcolor,
      linkcolor=linkcolor,
      citecolor=citecolor,
      }
    % Slightly bigger margins than the latex defaults
    
    \geometry{verbose,tmargin=1in,bmargin=1in,lmargin=1in,rmargin=1in}
    
    

\usepackage[brazil]{babel}
\usepackage{float}
\usepackage{url}
\usepackage{booktabs}
\usepackage[aboveskip=6pt,position=top]{caption}
\captionsetup[table]{position=top}
\captionsetup[figure]{position=bottom}
\usepackage{listings}
\usepackage{xcolor}
\definecolor{codegreen}{rgb}{0,0.6,0}
\definecolor{codegray}{rgb}{0.5,0.5,0.5}
\definecolor{codepurple}{rgb}{0.58,0,0.82}
\definecolor{backcolour}{rgb}{0.95,0.95,0.92}
\lstset{backgroundcolor=\color{backcolour},commentstyle=\color{codegreen},keywordstyle=\color{magenta},basicstyle=\ttfamily\small,breaklines=true,captionpos=t,keepspaces=true,numbers=left,numbersep=5pt,tabsize=4}
\renewcommand{\familydefault}{\sfdefault}
\renewcommand{\contentsname}{Sumário}
\renewcommand{\figurename}{Figura}
\renewcommand{\tablename}{Tabela}
\renewcommand{\refname}{Referências}
\begin{document}
    
    
    
    

    
    \begin{titlepage}
    \centering

    % 1. Logotipo do IF Goiano Local
    \includegraphics[width=0.15\textwidth]{logo_urutai.png}

    \vspace{1cm}

    % 2. Cabecalho Institucional
    {\small \textbf{INSTITUTO FEDERAL GOIANO} \\}
    {\small \textbf{CAMPUS URUTAÍ} \\}
    {\footnotesize NÚCLEO DE INFORMÁTICA \\}
    {\footnotesize ESPECIALIZAÇÃO EM INTELIGÊNCIA ARTIFICIAL APLICADA A DADOS CORPORATIVOS \\}
    \vspace{2cm}

    {\Large \MakeUppercase{Leonardo Maximino Bernardo} \par}
    \vspace{4.5cm}

    {\Large \textbf{\MakeUppercase{AcolheMente Escolar PB: Motor de Inferência Lógico para Apoio à Gestão Escolar na Priorização Responsável de Acolhimento em Saúde Mental}} \par}

    \vfill

    {\large Urutaí, maio de 2026 \par}
\end{titlepage}

    \begin{titlepage}
    \centering

    % 1. Cabecalho Institucional Simples
    {\small \textbf{INSTITUTO FEDERAL GOIANO} \\}
    {\small CAMPUS URUTAÍ \\}
    {\footnotesize Núcleo de Informática \\}
    {\footnotesize Pós-Graduação \textit{Lato Sensu} em Inteligência Artificial Aplicada a Dados Corporativos \\}

    \vspace{1.2cm}

    {\Large \textbf{Trabalho de Conclusão da Trilha de Aprendizagem I} \\}
    {\large \textit{Inteligência Artificial Simbólica e Busca} \\}

    \vspace{1.5cm}

    {\Large \textbf{AcolheMente Escolar PB} \\}
    \vspace{0.3cm}
    {\large Motor de Inferência Lógico para Apoio à Gestão Escolar \\
    na Priorização Responsável de Acolhimento em Saúde Mental}

    \vspace{1.5cm}

    {\large \textbf{Autor:} Leonardo Maximino Bernardo \\}
    {\normalsize \textbf{Orientação:} Corpo docente da Trilha I --- IF Goiano}

    \vfill

    {\large Urutaí --- GO \\}
    {\large Maio de 2026}
\end{titlepage}
\tableofcontents
\newpage


    \section{Apresentação do Aluno e
Contexto}\label{apresentauxe7uxe3o-do-aluno-e-contexto}

\begin{quote}
Leonardo Maximino Bernardo é professor de Física da rede pública
estadual da Paraíba desde 2020, com mais de 15 anos de experiência em
docência nos níveis médio e superior. Doutor em Ciência e Engenharia de
Materiais pela Universidade Federal da Paraíba, sua trajetória acadêmica
combina modelagem computacional avançada e uma crescente atuação em
ciência de dados e inteligência artificial.
\end{quote}

\begin{quote}
Ao longo de sua formação, atuou como pesquisador na UFPB, desenvolvendo
rotinas de cálculo e algoritmos para análise de grandes volumes de dados
de simulação numérica de solidificação de ligas metálicas, extraindo
padrões quantitativos para publicações acadêmicas. Essa experiência
consolidou habilidades de pensamento analítico, resolução de problemas
complexos e comunicação científica que informam diretamente sua prática
docente e sua abordagem em projetos de inteligência artificial.
\end{quote}

\begin{quote}
Atualmente, além da docência em Física, cursa Pós-Graduação \emph{Lato
Sensu} em Inteligência Artificial Aplicada pelo Instituto Federal Goiano
e possui formação em Ciência de Dados pelo Unipê. Sua \emph{stack}
tecnológica inclui Python, JavaScript, SQL, \emph{frameworks} como
Django e React, e ferramentas de análise como Pandas, Scikit-learn e
NetworkX. Possui certificações em desenvolvimento \emph{back-end} com
Node.js, Harvard CS50 e aceleração global de desenvolvimento de
software.
\end{quote}

\begin{quote}
O presente trabalho nasce da intersecção entre sua vivência como
professor de escola pública estadual na Paraíba --- onde testemunha
diariamente demandas crescentes de acolhimento em saúde mental de
adolescentes --- e seu domínio técnico em modelagem computacional e
lógica formal. A escolha do domínio de saúde mental escolar reflete o
compromisso de aplicar inteligência artificial de forma ética,
explicável e governada em um contexto de vulnerabilidade social, onde
erros tecnológicos podem ter consequências humanas graves.
\end{quote}

    \section{Introdução e Definição do
Problema}\label{introduuxe7uxe3o-e-definiuxe7uxe3o-do-problema}

\begin{quote}
A saúde mental de adolescentes no Brasil atravessa um momento crítico.
Dados da Pesquisa Nacional de Saúde do Escolar --- PeNSE 2024 (IBGE,
2024) --- revelam que parcela significativa dos estudantes de 13 a 17
anos reporta sentimentos persistentes de tristeza, solidão e
desesperança, com índices particularmente preocupantes na região
Nordeste e, especificamente, no estado da Paraíba. O relatório da
Organização Mundial da Saúde (WHO, 2022) confirma essa tendência global:
transtornos mentais respondem por uma proporção crescente da carga de
doença entre jovens, e a maioria dos casos permanece sem qualquer forma
de acolhimento adequado.
\end{quote}

\begin{quote}
No contexto da escola pública estadual paraibana, a demanda por
acolhimento em saúde mental frequentemente supera a capacidade de
resposta da equipe pedagógica e psicossocial. Professores, coordenadores
e orientadores educacionais enfrentam o desafio de identificar, dentre
centenas de estudantes, aqueles que necessitam de atenção prioritária
--- sem dispor de ferramentas estruturadas para essa tarefa. O resultado
é uma dependência excessiva de percepções individuais, sujeitas a viés
cognitivo, fadiga decisória e inconsistência entre turnos e unidades
escolares. Essa lacuna não é trivial: a ausência de priorização
sistemática pode significar que um adolescente em situação de risco
grave não receba acolhimento em tempo hábil.
\end{quote}

\begin{quote}
Ao mesmo tempo, a adoção acrítica de tecnologias de inteligência
artificial nesse domínio apresenta riscos igualmente sérios. Sistemas
baseados em aprendizado de máquina ou IA generativa, embora poderosos em
outros contextos, introduzem opacidade decisória incompatível com o
rigor ético exigido quando se trata de menores de idade em situação de
vulnerabilidade (Jobin; Ienca; Vayena, 2019). A Lei Geral de Proteção de
Dados --- LGPD (Brasil, 2018) --- e o Estatuto da Criança e do
Adolescente --- ECA (Brasil, 1990) --- impõem salvaguardas rigorosas
para o tratamento de dados sensíveis de menores, incluindo o direito à
explicação de decisões automatizadas. Um modelo de ``caixa preta'' que
classifique ou ranqueie estudantes por risco psicológico seria, no
mínimo, juridicamente temerário e, no limite, eticamente inaceitável.
\end{quote}

\begin{quote}
Diante desse cenário, o presente trabalho propõe o \textbf{AcolheMente
Escolar PB}: um Sistema Baseado em Conhecimento que utiliza Lógica
Proposicional com Inferência por Resolução para apoiar --- e nunca
substituir --- a gestão escolar na priorização responsável de
acolhimento em saúde mental. A escolha por lógica simbólica
determinística não é uma limitação técnica, mas uma decisão de
\emph{design} ético: garante explicabilidade total, rastreabilidade de
cada conclusão e ausência de qualquer linguagem diagnóstica clínica
(Russell; Norvig, 2022).
\end{quote}

\textbf{Problema central:} Como apoiar a gestão escolar de escolas
públicas estaduais da Paraíba na priorização sistemática de acolhimento
em saúde mental de adolescentes, utilizando inteligência artificial
explicável, sem produzir diagnósticos clínicos, sem violar a LGPD ou o
ECA, e sem substituir o julgamento humano profissional?

\textbf{Objetivo geral:} Desenvolver e validar um motor de inferência
por resolução, baseado em lógica proposicional e fundamentado em dados
oficiais da PeNSE 2024, para apoiar a priorização de acolhimento em
saúde mental escolar no estado da Paraíba.

\textbf{Objetivos específicos:} 1. Modelar o domínio em 7 variáveis
proposicionais (4 entradas nucleares, 2 entradas contextuais e 1 saída
inferida), ancoradas em variáveis da PeNSE 2024. 2. Formalizar 6 regras
de inferência em Forma Normal Conjuntiva e implementar motor de
resolução determinístico. 3. Validar o motor com cenários sintéticos
representativos e comparar com \emph{baseline} manual. 4. Garantir
conformidade integral com LGPD, ECA e Programa Saúde na Escola (Brasil,
2007). 5. Documentar explicabilidade, limitações e perspectivas de
evolução com rigor acadêmico.

    \section{Metodologia e Escolha da Técnica de
IA}\label{metodologia-e-escolha-da-tuxe9cnica-de-ia}

\begin{quote}
O AcolheMente Escolar PB é, em sua essência, um Sistema Baseado em
Conhecimento --- uma classe de sistemas de inteligência artificial que
codifica explicitamente o conhecimento de especialistas humanos em
regras formais e utiliza mecanismos de inferência para derivar
conclusões a partir de fatos observados (Brachman; Levesque, 2004).
Diferentemente de abordagens estatísticas ou conexionistas, um SBC não
``aprende'' a partir de dados brutos: opera sobre uma base de
conhecimento construída deliberadamente por engenheiros do conhecimento
em colaboração com especialistas do domínio.
\end{quote}

\begin{quote}
A técnica específica adotada é a \textbf{Lógica Proposicional com
Inferência por Resolução}. Trata-se de um formalismo da lógica clássica
no qual o conhecimento é representado por proposições atômicas e
conectivos lógicos. A inferência é realizada por meio do método de
resolução, um procedimento correto e refutacionalmente completo para a
lógica proposicional (Genesereth; Nilsson, 1987). O método opera sobre
cláusulas em Forma Normal Conjuntiva: para verificar se uma conclusão
\(A\) pode ser derivada de uma base de conhecimento \(KB\), adiciona-se
a negação da conclusão (\(\neg A\)) ao conjunto de cláusulas e aplica-se
repetidamente a regra de resolução, buscando derivar a cláusula vazia
--- o que constitui prova por contradição de que \(KB \models A\).
\end{quote}

\begin{quote}
Foram consideradas e descartadas alternativas como Lógica de Primeira
Ordem (complexidade desnecessária para 7 variáveis binárias), Redes
Bayesianas (exigiriam distribuições de probabilidade inestimáveis),
Árvores de Decisão (dependem de dados rotulados inexistentes), Redes
Neurais (incompatíveis com explicabilidade total) e IA Generativa (sem
garantias de determinismo).
\end{quote}

\begin{quote}
Conforme argumentam Russell e Norvig (2022), a escolha da representação
do conhecimento deve ser guiada pelas propriedades do domínio, e não
pela sofisticação da técnica. No caso do AcolheMente Escolar PB, a
lógica proposicional é a ferramenta \emph{correta}, não uma ferramenta
\emph{limitada}.
\end{quote}

\begin{quote}
O princípio de \emph{human-in-the-loop} é inegociável. A variável de
saída \(A\) é um sinal para o profissional humano, não uma decisão
final.
\end{quote}

\subsection{Arquitetura de Dados, Proveniência e
Governança}\label{arquitetura-de-dados-proveniuxeancia-e-governanuxe7a}

\begin{quote}
O AcolheMente opera com arquitetura rigorosa de proveniência de dados em
três camadas mutuamente exclusivas. A separação não é apenas técnica: é
um requisito de governança que impede contaminação cruzada entre fontes
de naturezas distintas.
\end{quote}

{\def\LTcaptype{none} % do not increment counter
\begin{longtable}[]{@{}
  >{\raggedright\arraybackslash}p{(\linewidth - 8\tabcolsep) * \real{0.1053}}
  >{\raggedright\arraybackslash}p{(\linewidth - 8\tabcolsep) * \real{0.1228}}
  >{\raggedright\arraybackslash}p{(\linewidth - 8\tabcolsep) * \real{0.2632}}
  >{\raggedright\arraybackslash}p{(\linewidth - 8\tabcolsep) * \real{0.2456}}
  >{\raggedright\arraybackslash}p{(\linewidth - 8\tabcolsep) * \real{0.2632}}@{}}
\toprule\noalign{}
\begin{minipage}[b]{\linewidth}\raggedright
Tier
\end{minipage} & \begin{minipage}[b]{\linewidth}\raggedright
Fonte
\end{minipage} & \begin{minipage}[b]{\linewidth}\raggedright
Uso Permitido
\end{minipage} & \begin{minipage}[b]{\linewidth}\raggedright
Uso Proibido
\end{minipage} & \begin{minipage}[b]{\linewidth}\raggedright
Justificativa
\end{minipage} \\
\midrule\noalign{}
\endhead
\bottomrule\noalign{}
\endlastfoot
TIER\_A & PeNSE 2024 (IBGE) & EDA, motivação, variáveis & Inferência
individual & Pesquisa populacional \\
TIER\_B & Cenários sintéticos & Validação do motor & Representar alunos
reais & Controle determinístico \\
TIER\_C & Dataset externo & Benchmark metodológico & Conclusão sobre
PB/BR & Sem validade regional \\
\end{longtable}
}

\begin{quote}
\textbf{Merge entre tiers é estritamente proibido.} Os dados da PeNSE
\textbf{não alimentam diretamente o motor de inferência} para
classificação individual: a PeNSE é pesquisa populacional, não sistema
de acompanhamento escolar ativo. O motor é validado por cenários
sintéticos determinísticos (TIER\_B), enquanto a PeNSE sustenta a
justificativa epidemiológica e a modelagem conceitual (TIER\_A).
\end{quote}

\begin{quote}
A LGPD (Brasil, 2018) classifica dados de saúde de menores como dados
pessoais sensíveis. O AcolheMente atende \emph{by design}: não processa
dados pessoais identificáveis, opera sobre indicadores agregados
anonimizados, toda conclusão é rastreável e não produz decisões
automatizadas sobre indivíduos. O ECA (Brasil, 1990) veda qualquer forma
de rotulagem de crianças e adolescentes.
\end{quote}

\subsection{Representação do Conhecimento: 7 Variáveis
Proposicionais}\label{representauxe7uxe3o-do-conhecimento-7-variuxe1veis-proposicionais}

\begin{quote}
A modelagem do domínio resultou em 7 variáveis proposicionais --- 4
entradas nucleares, 2 entradas contextuais e 1 saída inferida:
\end{quote}

{\def\LTcaptype{none} % do not increment counter
\begin{longtable}[]{@{}
  >{\raggedright\arraybackslash}p{(\linewidth - 6\tabcolsep) * \real{0.1042}}
  >{\raggedright\arraybackslash}p{(\linewidth - 6\tabcolsep) * \real{0.1250}}
  >{\raggedright\arraybackslash}p{(\linewidth - 6\tabcolsep) * \real{0.4792}}
  >{\raggedright\arraybackslash}p{(\linewidth - 6\tabcolsep) * \real{0.2917}}@{}}
\toprule\noalign{}
\begin{minipage}[b]{\linewidth}\raggedright
Var
\end{minipage} & \begin{minipage}[b]{\linewidth}\raggedright
Tipo
\end{minipage} & \begin{minipage}[b]{\linewidth}\raggedright
Semântica Operacional
\end{minipage} & \begin{minipage}[b]{\linewidth}\raggedright
Fontes PeNSE
\end{minipage} \\
\midrule\noalign{}
\endhead
\bottomrule\noalign{}
\endlastfoot
E & Nuclear & Sofrimento emocional recorrente & B12004, B12005,
B12007 \\
B & Nuclear & Baixo apoio socioafetivo percebido & B12003, B07004 \\
V & Nuclear & Indicador crítico de desvalor da vida & B12008 \\
S & Nuclear & Sinal autorreferido de autoagressão & B12009 \\
A & Nuclear & Acolhimento humano prioritário (saída) & Inferida \\
C & Contextual & Contexto comportamental agravante & B03010C, B03006B \\
I & Contextual & Insuficiência institucional de resposta & E01P60,
E01P117 \\
\end{longtable}
}

\begin{quote}
As variáveis contextuais \textbf{nunca inferem A isoladamente} --- este
é um \emph{guardrail} lógico fundamental.
\end{quote}

\subsection{Base de Conhecimento: 6 Regras
R1--R6}\label{base-de-conhecimento-6-regras-r1r6}

{\def\LTcaptype{none} % do not increment counter
\begin{longtable}[]{@{}
  >{\raggedright\arraybackslash}p{(\linewidth - 4\tabcolsep) * \real{0.2258}}
  >{\raggedright\arraybackslash}p{(\linewidth - 4\tabcolsep) * \real{0.6129}}
  >{\raggedright\arraybackslash}p{(\linewidth - 4\tabcolsep) * \real{0.1613}}@{}}
\toprule\noalign{}
\begin{minipage}[b]{\linewidth}\raggedright
Regra
\end{minipage} & \begin{minipage}[b]{\linewidth}\raggedright
Forma Implicativa
\end{minipage} & \begin{minipage}[b]{\linewidth}\raggedright
CNF
\end{minipage} \\
\midrule\noalign{}
\endhead
\bottomrule\noalign{}
\endlastfoot
R1 & \(S \rightarrow A\) & \(\neg S \lor A\) \\
R2 & \(V \land E \rightarrow A\) & \(\neg V \lor \neg E \lor A\) \\
R3 & \(E \land B \rightarrow A\) & \(\neg E \lor \neg B \lor A\) \\
R4 & \(V \land B \rightarrow A\) & \(\neg V \lor \neg B \lor A\) \\
R5 & \(E \land B \land C \rightarrow A\) &
\(\neg E \lor \neg B \lor \neg C \lor A\) \\
R6 & \(V \land I \rightarrow A\) & \(\neg V \lor \neg I \lor A\) \\
\end{longtable}
}

\begin{quote}
A conversão para CNF segue a equivalência
\(p \rightarrow q \equiv \neg p \lor q\) e sua generalização conjuntiva.
\end{quote}

    \section{Implementação Acadêmica e Motores
Equivalentes}\label{implementauxe7uxe3o-acaduxeamica-e-motores-equivalentes}

\begin{quote}
A implementação segue princípios de modularidade e testabilidade. A
arquitetura separa três responsabilidades: representação das cláusulas
CNF, mecanismo de resolução e interface de validação com cenários
sintéticos. O motor de inferência foi implementado em Python puro, sem
dependências externas.
\end{quote}

\begin{quote}
O projeto mantém duas versões do motor, ambas matematicamente
equivalentes. A versão de produção (\texttt{src/acolhemente/motor.py})
possui arquitetura modular com camadas de abstração e integração com o
grafo de explicabilidade. A versão acadêmica
(\texttt{appendices/motor\_resolucao\_acolhemente.py}), apresentada a
seguir, é uma simplificação didática com menos de 100 linhas, sem
dependências externas, com comentários explicativos em português. A
simplificação é didática, não matemática: as regras R1--R6, a CNF e a
variável de saída \(A\) são idênticas em ambas as versões.
\end{quote}

\begin{quote}
A seguir, cada apêndice é salvo via \texttt{\%\%writefile} e em seguida
executado para demonstração.
\end{quote}

    \begin{Verbatim}[commandchars=\\\{\}]
[Celula Colab — execucao reservada ao Google Colab]
    \end{Verbatim}

    \begin{Verbatim}[commandchars=\\\{\}]
[Celula Colab — execucao reservada ao Google Colab]
    \end{Verbatim}

    \begin{Verbatim}[commandchars=\\\{\}]
[Celula Colab — execucao reservada ao Google Colab]
    \end{Verbatim}

    \begin{Verbatim}[commandchars=\\\{\}]
======================================================================
REGRAS R1-R6: Forma Implicativa e CNF
======================================================================
  R1: S -> A                     CNF: \textasciitilde{}S OR A
  R2: V AND E -> A               CNF: \textasciitilde{}V OR \textasciitilde{}E OR A
  R3: E AND B -> A               CNF: \textasciitilde{}E OR \textasciitilde{}B OR A
  R4: V AND B -> A               CNF: \textasciitilde{}V OR \textasciitilde{}B OR A
  R5: E AND B AND C -> A         CNF: \textasciitilde{}E OR \textasciitilde{}B OR \textasciitilde{}C OR A
  R6: V AND I -> A               CNF: \textasciitilde{}V OR \textasciitilde{}I OR A

======================================================================
DEMONSTRACAO: Motor de Inferencia por Resolucao
======================================================================
Cenario                            A  Regras acionadas
----------------------------------------------------------------------
  Cenario Nulo                   NAO  --
  S isolado (R1)                 SIM  R1
  V+E (R2)                       SIM  R2
  E+B (R3)                       SIM  R3
  V+B (R4)                       SIM  R4
  E+B+C (R5)                     SIM  R3, R5
  V+I (R6)                       SIM  R6
  C isolado (guardrail)          NAO  --
  I isolado (guardrail)          NAO  --
  C+I (guardrail duplo)          NAO  --
  Todas verdadeiras              SIM  R1, R2, R3, R4, R5, R6

Motor de resolucao: CORRETO e REFUTACIONALMENTE COMPLETO
    \end{Verbatim}

    \begin{Verbatim}[commandchars=\\\{\}]
[Celula Colab — execucao reservada ao Google Colab]
    \end{Verbatim}

    \begin{Verbatim}[commandchars=\\\{\}]
================================================================================
VALIDACAO EXAUSTIVA: 14 Cenarios Sinteticos
================================================================================
Cenario                        Esperado   Obtido   Status  Regras
--------------------------------------------------------------------------------
  Cenario Nulo                      NAO      NAO       OK  --
  S isolado (R1)                    SIM      SIM       OK  R1
  V+E (R2)                          SIM      SIM       OK  R2
  E+B (R3)                          SIM      SIM       OK  R3
  V+B (R4)                          SIM      SIM       OK  R4
  E+B+C (R5)                        SIM      SIM       OK  R3, R5
  V+I (R6)                          SIM      SIM       OK  R6
  C isolado (guardrail)             NAO      NAO       OK  --
  I isolado (guardrail)             NAO      NAO       OK  --
  C+I (guardrail duplo)             NAO      NAO       OK  --
  Todas verdadeiras                 SIM      SIM       OK  R1, R2, R3, R4, R5,
R6
  E isolado                         NAO      NAO       OK  --
  B isolado                         NAO      NAO       OK  --
  V isolado                         NAO      NAO       OK  --

================================================================================
Total: 14 cenarios, 0 falha(s)

RESULTADO: TODOS OS CENARIOS PASSARAM
    \end{Verbatim}

    \subsubsection{Nota sobre R5 e subsunção
lógica}\label{nota-sobre-r5-e-subsunuxe7uxe3o-luxf3gica}

\begin{quote}
A regra R5 (\texttt{E\ ∧\ B\ ∧\ C\ →\ A}) é logicamente subsumida pela
regra R3 (\texttt{E\ ∧\ B\ →\ A}): sempre que R5 é verdadeira, R3 também
é verdadeira, pois \texttt{E=1} e \texttt{B=1} já bastam para R3.
Portanto, R5 nunca é a \emph{única} regra que aciona A. Mesmo assim, R5
é \textbf{mantida} na base de conhecimento por três razões: (1)
enriquece a explicação ao profissional humano, incluindo o fator
contextual C; (2) prepara a base para futuras versões com lógica fuzzy;
(3) não introduz falsos positivos. A validação exaustiva abaixo confirma
que em nenhum dos 64 cenários R5 altera o resultado.
\end{quote}

Referência:
\texttt{docs/adr/ADR-v3\_3-r5-subsuncao-e-validacao-exaustiva.md}

    \subsubsection{Nota metodológica de
governança}\label{nota-metodoluxf3gica-de-governanuxe7a}

A PeNSE é utilizada como referência epidemiológica e ancoragem
conceitual para escolha das variáveis proposicionais, não como entrada
operacional do motor de inferência individual.

A variável A é saída inferida, não variável nuclear observada.

A regra R5 é logicamente subsumida por R3: quando E, B e C são
verdadeiros, E e B já bastam para acionar R3. R5 é mantida como marcador
explicativo contextual, sem alterar o resultado decisório.

A comparação entre baseline manual e motor lógico é uma rubrica
qualitativa, não experimento empírico com profissionais da escola.

    \begin{Verbatim}[commandchars=\\\{\}]
Validação exaustiva 2\^{}6 = 64 combinações
Total testado: 64
Falhas: 0
A=SIM: 50
A=NÃO: 14
R5 sem R3: 0
CSV salvo em: validacao\_exaustiva\_64.csv
Status: APROVADO
    \end{Verbatim}

    \begin{Verbatim}[commandchars=\\\{\}]
[Celula Colab — execucao reservada ao Google Colab]
    \end{Verbatim}

    \begin{Verbatim}[commandchars=\\\{\}]
[Celula Colab — execucao reservada ao Google Colab]
    \end{Verbatim}

    \section{Resultados e Comparações}\label{resultados-e-comparauxe7uxf5es}

\begin{quote}
O motor foi validado com o conjunto exaustivo de \(2^6 = 64\)
combinações possíveis das variáveis de entrada. Para cada cenário,
verificou-se se a conclusão do motor está correta em relação às regras
formais, quais regras foram acionadas e se a rastreabilidade é completa.
\end{quote}

\begin{quote}
Os resultados confirmam o comportamento esperado em todos os cenários
testados. Quando nenhuma variável é verdadeira (cenário nulo), \(A\) não
é inferida, demonstrando ausência de falsos positivos. A variável \(S\)
(autoagressão) é suficiente para acionar o acolhimento (R1). As
variáveis contextuais (\(C\), \(I\), \(C+I\)), quando verdadeiras
isoladamente, \textbf{não acionam} \(A\), confirmando o \emph{guardrail}
arquitetural.
\end{quote}

{\def\LTcaptype{none} % do not increment counter
\begin{longtable}[]{@{}lllllllll@{}}
\toprule\noalign{}
Cenário & E & B & V & S & C & I & A & Regra(s) \\
\midrule\noalign{}
\endhead
\bottomrule\noalign{}
\endlastfoot
Nulo & 0 & 0 & 0 & 0 & 0 & 0 & NÃO & --- \\
S isolado & 0 & 0 & 0 & 1 & 0 & 0 & SIM & R1 \\
V+E & 1 & 0 & 1 & 0 & 0 & 0 & SIM & R2 \\
E+B & 1 & 1 & 0 & 0 & 0 & 0 & SIM & R3 \\
V+B & 0 & 1 & 1 & 0 & 0 & 0 & SIM & R4 \\
E+B+C & 1 & 1 & 0 & 0 & 1 & 0 & SIM & R3, R5 \\
V+I & 0 & 0 & 1 & 0 & 0 & 1 & SIM & R6 \\
C isolado & 0 & 0 & 0 & 0 & 1 & 0 & NÃO & --- \\
I isolado & 0 & 0 & 0 & 0 & 0 & 1 & NÃO & --- \\
Todas & 1 & 1 & 1 & 1 & 1 & 1 & SIM & R1--R6 \\
\end{longtable}
}

\subsubsection{\texorpdfstring{Comparação: \emph{Baseline} Manual
vs.~Motor
Lógico}{Comparação: Baseline Manual vs.~Motor Lógico}}\label{comparauxe7uxe3o-baseline-manual-vs.-motor-luxf3gico}

{\def\LTcaptype{none} % do not increment counter
\begin{longtable}[]{@{}lll@{}}
\toprule\noalign{}
Critério & Manual & Motor Lógico \\
\midrule\noalign{}
\endhead
\bottomrule\noalign{}
\endlastfoot
Transparência & Baixa & Total \\
Consistência & Variável & Determinística \\
Explicabilidade & Verbal & Formal (CNF) \\
Rastreabilidade & Nenhuma & Regra por regra \\
Fadiga decisória & Presente & Ausente \\
Nuance qualitativa & Alta & Nenhuma \\
\end{longtable}
}

\subsection{Explicabilidade}\label{explicabilidade}

\begin{quote}
A explicabilidade é um requisito ético, jurídico e pedagógico
simultaneamente. O AcolheMente implementa explicabilidade
\textbf{intrínseca}: a lógica proposicional é, por definição,
transparente. Cada variável tem semântica clara, cada regra tem
antecedentes e consequente explícitos, e cada conclusão pode ser
rastreada à cláusula CNF que a originou.
\end{quote}

\begin{quote}
Ribeiro, Singh e Guestrin (2016) definem explicabilidade como a
capacidade de um sistema de apresentar, em termos compreensíveis, as
razões que levaram a uma determinada conclusão. No AcolheMente, a
explicação \emph{é} o modelo --- não há gap entre o que o sistema faz e
o que o sistema diz que faz.
\end{quote}

\begin{quote}
A distinção entre explicabilidade intrínseca e \emph{post-hoc} é
particularmente relevante no domínio de saúde mental escolar. O ECA
(Brasil, 1990) exige proteção integral de crianças e adolescentes, o que
inclui proteção contra classificação automatizada opaca. A LGPD (Brasil,
2018) garante transparência no tratamento de dados.
\end{quote}

\subsection{Discussão Crítica}\label{discussuxe3o-cruxedtica}

\begin{quote}
O AcolheMente demonstra que é possível aplicar inteligência artificial
em um domínio sensível com explicabilidade total, determinismo e
conformidade jurídica. A escolha por lógica proposicional é precisamente
adequada ao escopo do problema.
\end{quote}

\begin{quote}
\textbf{Limitações reconhecidas:} A lógica proposicional não modela
gradações, incertezas ou relações temporais. A qualidade das conclusões
depende da alimentação correta das variáveis de entrada. O motor não
aprende com novos dados. A validação foi computacional, não em ambiente
escolar real.
\end{quote}

\begin{quote}
\textbf{Riscos:} Falso conforto tecnológico (equipe escolar reduz
atenção qualitativa), viés de construção (regras refletem julgamento dos
engenheiros), risco de medicalização do espaço escolar.
\end{quote}

\begin{quote}
A contribuição do projeto reside na demonstração de que rigor técnico e
responsabilidade ética podem coexistir --- e de que, em domínios de alto
risco humano, a transparência é mais valiosa que a complexidade.
\end{quote}

    \section{Perspectiva de Evolução do
Trabalho}\label{perspectiva-de-evoluuxe7uxe3o-do-trabalho}

\begin{quote}
A arquitetura do AcolheMente Escolar PB foi concebida para permitir
evolução incremental, mas essa evolução não é governada pela lógica do
``mais complexo é melhor''. Qualquer avanço técnico deve ser precedido
por validações éticas e institucionais proporcionais ao risco que a
mudança introduz.
\end{quote}

\begin{quote}
O caminho mais natural de evolução é o refinamento do próprio motor
simbólico, sem abandonar o paradigma determinístico. A incorporação de
novas variáveis --- frequência escolar, participação em atividades
extracurriculares, indicadores de convivência familiar --- poderia
enriquecer a modelagem sem comprometer a transparência. Essa expansão,
entretanto, não é trivial: cada nova variável dobra o espaço
combinatório de estados, exigindo revisão integral das regras e
revalidação dos cenários.
\end{quote}

\begin{quote}
Uma evolução particularmente promissora é a transição para lógica
\emph{fuzzy}, com graus de pertinência entre 0 e 1, permitindo modelar
gradações de sofrimento emocional e apoio social. A definição das
funções de pertinência exigiria calibração empírica com dados
longitudinais e validação com profissionais de psicologia escolar
(Russell; Norvig, 2022).
\end{quote}

\begin{quote}
Em horizonte mais distante, modelos híbridos poderiam integrar módulos
de \emph{machine learning} para tarefas específicas, mantendo o motor
simbólico como camada de governança e explicabilidade. Nessa
arquitetura, nenhuma conclusão do módulo estatístico seria apresentada à
equipe escolar sem passar pelo crivo das regras lógicas (Jobin; Ienca;
Vayena, 2019).
\end{quote}

\begin{quote}
A utilização de redes neurais profundas permanece uma hipótese
estritamente teórica, condicionada a aprovação por comitê de ética,
conformidade demonstrável com LGPD (Brasil, 2018) e ECA (Brasil, 1990),
e base longitudinal com volume e qualidade suficientes.
\end{quote}

\begin{quote}
Qualquer evolução deve ser acompanhada de piloto institucional, formação
continuada, comitê de governança de IA e avaliação de impacto. O
\emph{roadmap} prevê integração com os fluxos do Programa Saúde na
Escola (Brasil, 2007).
\end{quote}

\begin{quote}
Versões futuras devem implementar monitoramento contínuo de viés e
auditoria humana periódica das regras por comitê multidisciplinar. O
risco de medicalização da escola deve ser permanentemente monitorado.
\end{quote}

\begin{quote}
É necessário resistir à tentação de equiparar sofisticação algorítmica a
melhoria do projeto. A simplicidade do motor proposicional é uma
\emph{feature} a ser preservada enquanto as condições éticas, jurídicas
e institucionais para abordagens mais complexas não estiverem plenamente
atendidas.
\end{quote}

\begin{quote}
Independentemente da evolução técnica, o compromisso fundamental
permanece inalterável: o sistema nunca diagnosticará condições clínicas.
Esse compromisso não é uma restrição técnica: é a razão de ser do
projeto.
\end{quote}

    \section{Conclusões}\label{conclusuxf5es}

\begin{quote}
O presente trabalho demonstrou a viabilidade de aplicar inteligência
artificial simbólica --- especificamente, lógica proposicional com
inferência por resolução --- ao domínio sensível de saúde mental
escolar. O AcolheMente Escolar PB, motor de inferência com 7 variáveis
proposicionais e 6 regras em Forma Normal Conjuntiva, oferece uma
contribuição técnica e metodológica ao debate sobre IA responsável em
contextos educacionais.
\end{quote}

\begin{quote}
\textbf{Contribuição técnica:} Formalização de domínio complexo em
lógica proposicional, demonstrando que problemas reais de apoio à
decisão podem ser modelados com técnicas simbólicas; motor de resolução
determinístico validado contra todos os \(2^6\) cenários possíveis;
\emph{pipeline} reprodutível com testes automatizados.
\end{quote}

\begin{quote}
\textbf{Contribuição metodológica:} Arquitetura de governança de dados
em três \emph{tiers}, \emph{guardrails} de não diagnóstico,
\emph{human-in-the-loop} inegociável, conformidade demonstrável com
LGPD, ECA e PSE.
\end{quote}

\begin{quote}
\textbf{Contribuição escolar:} Ferramenta de apoio à gestão que reduz a
dependência de percepções individuais, com mecanismo de explicabilidade
que permite auditar cada decisão.
\end{quote}

\begin{quote}
Os limites são reconhecidos: expressividade restrita, dependência de
alimentação humana correta, validação limitada a cenários sintéticos e
risco de falso conforto tecnológico. Esses limites não são falhas de
\emph{design}: são consequências deliberadas de um sistema que prioriza
segurança, transparência e cautela.
\end{quote}

\begin{quote}
A maturidade de um projeto de IA aplicada a populações vulneráveis não
se mede pela complexidade do modelo, mas pela robustez de sua
governança. O AcolheMente Escolar PB demonstra que é possível ser
tecnicamente rigoroso e eticamente responsável --- e que, em domínios de
alto risco humano, essa combinação não é opcional.
\end{quote}

    \section{Referências
Bibliográficas}\label{referuxeancias-bibliogruxe1ficas}

\begin{quote}
BRACHMAN, R. J.; LEVESQUE, H. J. \textbf{Knowledge Representation and
Reasoning}. San Francisco: Morgan Kaufmann, 2004.
\end{quote}

\begin{quote}
BRASIL. Lei nº 8.069, de 13 de julho de 1990 --- Estatuto da Criança e
do Adolescente (ECA). Disponível em:
https://www.planalto.gov.br/ccivil\_03/leis/L8069.htm. Acesso em: 20
maio 2026.
\end{quote}

\begin{quote}
BRASIL. Decreto nº 6.286, de 5 de dezembro de 2007 --- Programa Saúde na
Escola (PSE). Disponível em:
https://www.planalto.gov.br/ccivil\_03/\_ato2007-2010/2007/decreto/d6286.htm.
Acesso em: 20 maio 2026.
\end{quote}

\begin{quote}
BRASIL. Lei nº 13.709, de 14 de agosto de 2018 --- Lei Geral de Proteção
de Dados Pessoais (LGPD). Disponível em:
https://www.planalto.gov.br/ccivil\_03/\_ato2015-2018/2018/lei/L13709.htm.
Acesso em: 20 maio 2026.
\end{quote}

\begin{quote}
GENESERETH, M. R.; NILSSON, N. J. \textbf{Logical Foundations of
Artificial Intelligence}. San Mateo: Morgan Kaufmann, 1987.
\end{quote}

\begin{quote}
INSTITUTO BRASILEIRO DE GEOGRAFIA E ESTATÍSTICA. \textbf{Pesquisa
Nacional de Saúde do Escolar --- PeNSE 2024}. Rio de Janeiro: IBGE,
2024. Disponível em:
https://www.ibge.gov.br/estatisticas/sociais/educacao/9134-pesquisa-nacional-de-saude-do-escolar.html.
Acesso em: 20 maio 2026.
\end{quote}

\begin{quote}
JOBIN, A.; IENCA, M.; VAYENA, E. The global landscape of AI ethics
guidelines. \textbf{Nature Machine Intelligence}, v. 1, p.~389--399,
2019.
\end{quote}

\begin{quote}
RIBEIRO, M. T.; SINGH, S.; GUESTRIN, C. ``Why Should I Trust You?'':
Explaining the Predictions of Any Classifier. In: \textbf{Proceedings of
the 22nd ACM SIGKDD International Conference on Knowledge Discovery and
Data Mining}, p.~1135--1144. ACM, 2016.
\end{quote}

\begin{quote}
RUSSELL, S.; NORVIG, P. \textbf{Inteligência Artificial: uma abordagem
moderna}. 4. ed.~Rio de Janeiro: GEN LTC, 2022.
\end{quote}

\begin{quote}
WORLD HEALTH ORGANIZATION. \textbf{World mental health report:
transforming mental health for all}. Geneva: WHO, 2022. Disponível em:
https://www.who.int/publications/i/item/9789240049338. Acesso em: 20
maio 2026.
\end{quote}

    \newpage
\makeatletter
\renewcommand{\@chapapp}{Apêndice}
\makeatother

\clearpage
\appendix

\section*{Apêndice A --- Símbolos Proposicionais do AcolheMente}
\addcontentsline{toc}{section}{Apêndice A --- Símbolos Proposicionais do AcolheMente}
\lstinputlisting[language=Python, caption={Símbolos Proposicionais do AcolheMente.}, label={lst:simbolos}]{simbolos_acolhemente.py}

\section*{Apêndice B --- Regras Lógicas e CNF}
\addcontentsline{toc}{section}{Apêndice B --- Regras Lógicas e CNF}
\lstinputlisting[language=Python, caption={Regras Lógicas e CNF.}, label={lst:regras}]{regras_acolhemente.py}

\section*{Apêndice C --- Motor de Inferência por Resolução}
\addcontentsline{toc}{section}{Apêndice C --- Motor de Inferência por Resolução}
\lstinputlisting[language=Python, caption={Motor de Inferência por Resolução.}, label={lst:motor}]{motor_resolucao_acolhemente.py}

\section*{Apêndice D --- Cenários Sintéticos de Validação}
\addcontentsline{toc}{section}{Apêndice D --- Cenários Sintéticos de Validação}
\lstinputlisting[language=Python, caption={Cenários Sintéticos de Validação.}, label={lst:cenarios}]{cenarios_sinteticos_acolhemente.py}

\section*{Apêndice E --- Grafo de Explicabilidade}
\addcontentsline{toc}{section}{Apêndice E --- Grafo de Explicabilidade}
\lstinputlisting[language=Python, caption={Grafo de Explicabilidade.}, label={lst:grafo}]{grafo_explicabilidade_acolhemente.py}

\section*{Apêndice F --- Gerador de Tabelas de Resultados}
\addcontentsline{toc}{section}{Apêndice F --- Gerador de Tabelas de Resultados}
\lstinputlisting[language=Python, caption={Gerador de Tabelas de Resultados.}, label={lst:tabelas}]{gerar_tabelas_resultados.py}


\end{document}